\color{unimain}{\subsection{Токовая защита нулевой последовательности}}
\color{black}

\begin{enumerate}
%%%%%%%%%%%%%%%%%%%%%%%%%%%%%%%%%%%%%%%%%%%%%%%%%% ОБЩИЕ СВЕДЕНИЯ %%%%%%%%%%%%%%%%%%%%%%%%%%%%%%%%%%%%%%%%%%%%%%%%%%
\item Общие сведения

\begin{enumerate}
\item
Токовая защита нулевой последовательности (ТЗНП) предназначена для селективного
отключения поврежденного участка электрической сети при коротких замыканиях на землю, реагируя на
повышение тока нулевой последовательности.

\item
В составе логики токовой направленной защиты нулевой последовательности предусмотрено четыре
ступени от замыканий на землю.

\item
На рисунке \ref{pict_tnznp:TNZNP} приведена функциональная схема ТЗНП-1, выполнение функциональных схем
ступеней ТЗНП-2(3,4) аналогичное.

\item
Уставочные параметры функции ТЗНП представлены в таблице~\ref{app_set:tznp}.

\begin{figure}[H]
\centering
\includegraphics[width=0.95\textwidth,height=0.95\textheight,keepaspectratio]{\fbpath/08.01 СВ/661.0121 ТЗНП/img/TNZNP.pdf}
\captionof{figure}{Функциональная схема ТЗНП-1(2…4)}\label{pict_tnznp:TNZNP}
\end{figure}

\item
Пороги тока срабатывания ступеней ТЗНП задаются уставками \textbf{«3I0ср ТЗНП-1(2…4)»}. Действие
ступеней осуществляется с регулируемыми независимыми выдержками времени \textbf{«Тср ТЗНП-1(2…4)»}. Ввод
в работу ступеней ТЗНП осуществляется с помощью программных ключей \textbf{«ТЗНП-1(2…4)»}. 


\item
Имеется возможность задать различные режимы работы ступеней ТЗНП при возникновении
неисправности в цепях напряжения. Вариант действия логики ступени ТЗНП при
срабатывании БНН задается программным переключателем \textbf{«БНН ТЗНП-1(2…4)»}, имеющим три положения:
\begin{itemize}
	\item \textit{«Выведено»} -- режим работы ступени не зависит от БНН;
	\item \textit{«Вывод напр.»} -- ступень переводится в ненаправленный режим работы;
	\item \textit{«Вывод ст.»} --   при срабатывании БНН ступень выводится из работы (однако, если ступень задана ненаправленной, то она останется в работе независимо от срабатывания БНН; ступень также не будет выводиться из работы, если направленность автоматически вывелась при срабатывании ТЗНП).
\end{itemize}

\item
Предусмотрен автоматический вывод направленности всех ступеней ТЗНП при срабатывании ТЗНП,
это обеспечивает устойчивое срабатывание ТЗНП при неполнофазном отключении, что необходимо для
надежного действия УРОВ. Данный режим вводится с помощью программного переключателя \textbf{«Вывод
напр при ср ТЗНП»}.

\item
Для обеспечения несрабатывания ТЗНП от бросков намагничивающего тока предусмотрена
блокировка ступеней ТЗНП по содержанию второй гармоники в токе нулевой последовательности.
Блокировка основана на контроле отношения действующего значения второй гармоники тока нулевой
последовательности к основной гармонике. Ступени ТЗНП блокируется при превышении данного
соотношения пороговой величины, заданной уставкой \textbf{«3I02г/3I01г ТЗНП»}. Наличие функции блокировки от
бросков намагничивающего тока задается для каждой ступени отдельно с помощью программных ключей
\textbf{«БНТ ТЗНП-1(2…4)»}. Для предотвращения излишней блокировки при больших токах КЗ,
сопровождающихся насыщением ТТ, блокировка выводится из действия при превышении тока 3I0 основной
гармоники выше уставки \textbf{«3I0блок БНТ ТЗНП»}.

\item
Функциональная схема логики органа выявления БНТ приведена на рисунке~\ref{pict_tnznp:BNT}.

\begin{figure}[H]
\centering
\includegraphics[width=0.7\textwidth,height=0.7\textheight,keepaspectratio]{\fbpath/08.01 СВ/661.0121 ТЗНП/img/BNT.pdf}
\captionof{figure}{ Функциональная схема логики органа выявления БНТ}\label{pict_tnznp:BNT}
\end{figure}

\item
Оперативный ввод/вывод чувствительных ступеней ТЗНП производится с помощью команды \textit{«ТЗНП: Блокировка ЧС»},
а выбор выводимых ступеней ТЗНП определяется с помощью программных ключей \textbf{«Вывод ТЗНП-1(2…4) по чувств»}.

\item
ТЗНП может быть оперативно выведена из работы от внешнего сигнала \textbf{«ТЗНП: Вывод»}, назначенного на соответствующий дискретный вход или функциональную клавишу.
Выводятся все ступени ТЗНП, а также все виды ускорения ТЗНП. 

\end{enumerate}


%%%%%%%%%%%%%%%%%%%%%%%%%%%%%%%%%%%%%%%%%%%%%%%%%% ОНМ %%%%%%%%%%%%%%%%%%%%%%%%%%%%%%%%%%%%%%%%%%%%%%%%%%
\item Органы направления мощности

\begin{enumerate}
\item
Направленность ТЗНП обеспечивается органами направления мощности нулевой последовательности
(ОНМНП) и обратной последовательности (ОНМОП).

\item
ОНМНП реагирует на величины векторов тока и напряжения нулевой последовательности от ТН 1СШ, а
также угол сдвига между ними.

\item
Функциональная схема логики ОНМНП приведена на рисунке~\ref{pict_tnznp:onmnp}.

\begin{figure}[H]
\centering
\includegraphics[width=0.7\textwidth,height=0.7\textheight,keepaspectratio]{\fbpath/08.01 СВ/661.0121 ТЗНП/img/ONM_NP.pdf}
\captionof{figure}{Функциональная схема логики ОНМНП}\label{pict_tnznp:onmnp}
\end{figure}

\item
Так как мощность нулевой последовательности всегда направлена от места несимметрии во внешнюю
сеть, то ОНМНП в 1СШ срабатывает при направлении мощности нулевой последовательности от 1СШ к ШСВ
(при КЗ впереди), а  ОНМНП в сеть – при обратном направлении мощности. 

\item
За положительное направление отсчета векторов токов и напряжений принято направление против
часовой стрелки. Углы максимальной чувствительности откладываются от вектора напряжения 3U0 (вектор
поляризации) и задаются уставками \textbf{«Фмч ОНМНП-1СШ»} и \textbf{«Фмч ОНМНП-2СШ»}. Уставки по углам
максимальной чувствительности рекомендуется задавать для ОНМНП в 1СШ – плюс 110º, а для ОНМНП в
2СШ – плюс 290º. Ширина зоны срабатывания ± 85º от угла максимальной чувствительности.

\item
Пороги срабатывания ОНМНП в 1СШ и ОНМНП в 2СШ по току и напряжению регулируются уставками 
\textbf{«3I0ср ОНМНП-1СШ»}, \textbf{«3U0ср ОНМНП-1СШ»}, \textbf{«3I0ср ОНМНП-2СШ»} и \textbf{«3U0ср ОНМНП-2СШ»}. Выбор
напряжения нулевой последовательности, по которому работает ОНМНП, осуществляется с помощью
программного переключателя \textbf{«Выбор 3U0»}. 

\item
Орган направления мощности обратной последовательности ОНМОП в 1СШ срабатывает при
направлении мощности обратной последовательности от 1СШ к ШСВ, а ОНМОП в 2СШ -- при обратном
направлении мощности. 

\item
Функциональная схема логики ОНМОП приведена на рисунке~\ref{pict_tnznp:onmop}.

\begin{figure}[H]
\centering
\includegraphics[width=0.7\textwidth,height=0.7\textheight,keepaspectratio]{\fbpath/08.01 СВ/661.0121 ТЗНП/img/ONM_OP.pdf}
\captionof{figure}{Функциональная схема логики ОНМОП}\label{pict_tnznp:onmop}
\end{figure}

\item
Углы максимальной чувствительности задаются уставками \textbf{«Фмч ОНМОП-1СШ»} и \textbf{«Фмч ОНМОП-
2СШ»}. Уставки по углам максимальной чувствительности рекомендуется задавать для ОНМОП в 1СШ -- плюс
$110^\circ$, а для ОНМОП в 2СШ -- плюс $290^\circ$. Ширина зоны срабатывания $\pm 85^\circ$ от угла максимальной
чувствительности.

\item
Пороги срабатывания ОНМОП в 1СШ и ОНМОП в 2СШ  по току и напряжению регулируются уставками
\textbf{«I2ср ОНМОП-1СШ»}, \textbf{«U2ср ОНМОП-1СШ»}, \textbf{«I2ср ОНМОП-2СШ»}, и \textbf{«U2ср ОНМОП-2СШ»}.

\small
\setlength{\extrarowheight}{0.05cm} % высота ячеек в таблице
\begin{longtable}
{|>{\centering\arraybackslash}m{15.2cm}|>{\centering\arraybackslash}m{1.7cm}|}
\caption{Параметры ОНМНП(ОНМОП)\hfill\vspace{-0.5\baselineskip}}\label{tnznp:tbl1}\\ 
\hhline{|-|-|}
\rowcolor{gray!30}
Наименование параметра & Значение  \\ 
\hhline{|-|-|}
\endfirsthead

\caption*{\hspace{3pt}\emph{Продолжение таблицы \ref{tnznp:tbl1}\hfill\vspace{-0.5\baselineskip}}} \\
\hhline{|-|-|}
\rowcolor{gray!30}
Наименование параметра & Значение  \\ 
\hhline{|-|-|}
\endhead
\hhline{|-|-|}

\endfoot
\endlastfoot

\raggedright Время срабатывания ОНМ при утроенных по отношению к порогам срабатывания значениях тока и напряжения, мс, не более & \centering\arraybackslash 35  \\  \hline
\raggedright Время возврата ОНМ при одновременном сбросе входных величин тока и напряжения от номинальных значений до нуля, мс, не более  & \centering\arraybackslash 25  \\  \hline
\raggedright Средняя основная погрешность порогов срабатывания ОНМ по току и напряжению, \% уставки, не более   & \centering\arraybackslash $\pm$5  \\  \hline
\raggedright Средняя основная абсолютная погрешность ОНМ по углу максимальной чувствительности, эл. град., не более & \centering\arraybackslash $\pm$5 \\  \hline
\raggedright Дополнительная погрешность параметров срабатывания от изменения температуры окружающего воздуха в рабочем диапазоне, \% от среднего значения параметров, измеренных при температуре $(20\pm5)^\circ$C, не более & \centering\arraybackslash $\pm$5 \\  \hline
\raggedright Минимальная угловая ширина зоны срабатывания ОНМ, эл. град, не менее & \centering\arraybackslash 165 \\ \hline
\raggedright Коэффициент возврата ОНМ НП по напряжению и току, не менее  & \centering\arraybackslash 0,9\\  \hline

\end{longtable}
\setlength{\extrarowheight}{0.0cm} % возврат параметра высота ячеек в таблице
\normalsize

\item
На рисунке \ref{pict_tnznp:ONMNP} рисунке представлены характеристики срабатывания ОНМНП в 1СШ и ОНМНП в 2СШ. Аналогичные
характеристики для ОНМОП в 1СШ и ОНМОП в 2СШ.

\medskip

\begin{figure}[H]
\centering
\begin{subfigure}{0.5\textwidth}
    \centering
    \includegraphics[width=\textwidth, trim=0 0mm 125mm 0, clip]{\fbpath/08.01 СВ/661.0121 ТЗНП/img/ONMNP.pdf}
	    \caption{a) ОНМНП в 1СШ}
    \label{pict_tnznp:ONMNP_a}
\end{subfigure}
%\hfill
\begin{subfigure}{0.49\textwidth}
    \centering
    \includegraphics[width=\textwidth, trim=125mm 0mm 0 0, clip]{\fbpath/08.01 СВ/661.0121 ТЗНП/img/ONMNP.pdf}
    \caption{б) ОНМНП в 2СШ}
    \label{pict_tnznp:ONMNP_b}
\end{subfigure}

\caption{Характеристики срабатывания ОНМНП}
\label{pict_tnznp:ONMNP}
\end{figure}

\medskip

\item
Для выбора режима направленности по мощности нулевой последовательности для каждой ступени
ТЗНП предусмотрен программный переключатель \textbf{«Направл ТЗНП-1(2…4)»}, имеющий три положения:
\begin{itemize}
    \item \textit{«Ненапр»} -- ступень выполнена ненаправленной;
    \item \textit{«В 1СШ»} -- ступень выполнена прямонаправленной в сторону первой системы шин (1СШ);
    \item \textit{«В 2СШ»} -- ступень выполнена обратнонаправленной в сторону второй системы шин (2СШ).
\end{itemize}

\item
Кроме того, можно установить различные алгоритмы контроля направленности с помощью программных
переключателей \textbf{«Контр Направл ТЗНП-1(2…4)»}, имеющих три положения:
\begin{itemize}
    \item \textit{«Разр»} --  прямая направленность обеспечивается срабатыванием ОНМНП в 1СШ, обратная -- срабатыванием ОНМНП в 2СШ;
    \item \textit{«Блок»} -- прямая направленность обеспечивается отсутствием срабатывания ОНМНП в 2СШ, обратная -- отсутствием срабатывания ОНМНП в 1СШ;
    \item \textit{«Разр или блок»} -- прямая направленность обеспечивается срабатыванием ОНМНП в 1СШ или отсутствием срабатывания ОНМНП в 2СШ, обратная -- срабатыванием ОНМНП в 2СШ или отсутствием срабатывания ОНМНП в 1СШ.
\end{itemize}

\item
Положение переключателя \textit{«Разр или блок»} рекомендуется выставлять для ступеней ТЗНП,
выполняющих функцию дальнего резервирования, так как при удаленных КЗ напряжение $3U_0$ может оказаться
недостаточным для срабатывания ОНМ из-за его уменьшения по мере удаления от точки КЗ.

\item
Уставки для выбора режима направленности по мощности обратной последовательности аналогичные.

\item
Тип ОНМ, используемый для определения направленности ступени ТЗНП, выбирается программным
переключателем \textit{«Тип ОНМ ТЗНП-1(2…4)»}, имеющим четыре положения:
\begin{itemize}
    \item \textit{«НП»} -- для определения направленности ступени учитывается только ОНМНП;
    \item \textit{«ОП»} -- для определения направленности ступени учитывается только ОНМОП;
    \item \textit{«НП или ОП»} -- для определения направленности ступени учитывается один из органов, ОНМНП или ОНМОП;
    \item \textit{«НП и ОП»} -- для определения направленности ступени одновременно учитываются оба органа, ОНМНП и ОНМОП.
\end{itemize}

\end{enumerate}


%%%%%%%%%%%%%%%%%%%%%%%%%%%%%%%%%%%%%%%%%%%%%%%%%% ОУ %%%%%%%%%%%%%%%%%%%%%%%%%%%%%%%%%%%%%%%%%%%%%%%%%%
\item Оперативное ускорение (ОУ)

\begin{enumerate}

\item
В составе логики токовой направленной защиты нулевой последовательности предусмотрено два
режима оперативного ускорения ТЗНП.

\item
Функциональная схема ОУ1 ТЗНП приведена на рисунке~\ref{pict_tnznp:OA_TNZNP}.

\vspace{2mm}
\begin{figure}[H]
\centering
\includegraphics[width=0.9\textwidth,height=0.9\textheight,keepaspectratio]{\fbpath/08.01 СВ/661.0121 ТЗНП/img/OA_TNZNP.pdf}
\captionof{figure}{Функциональная схема ОУ1 ТЗНП}\label{pict_tnznp:OA_TNZNP}
\end{figure}

\item
Оперативное ускорение ТЗНП с первой выдержкой времени предназначено для оперативного
ускорения ступени ТЗНП при выводе ДЗШ. Ввод в работу ОУ ТЗНП и выбор ускоряемой ступени задается
программным переключателем \textbf{«Ступень ОУ1 ТЗНП»}. Предусмотрена возможность вывода
направленности ускоряемой ступени с помощью программного ключа \textbf{«Вывод напр при ср ТЗНП»}.
Действие ускорения осуществляется с регулируемой независимой выдержкой времени \textbf{«Тср ОУ1 ТЗНП»}. При
выводе направленности ускоряемой ступени ее блокировка от БНТ осуществляется с помощью программного
ключа \textbf{«БНТ ОУ1 ТЗНП»}.

\item
Оперативный ввод ОУ1 ТЗНП производится с помощью команды управления \textit{«ТЗНП/ОУ1 ТЗНП: Ввод»}.

\item
Функциональная схема ОУ2 ТЗНП приведена на рисунке~\ref{pict_tnznp:OA_TNZNP_without_T}.

\vspace{2mm}
\begin{figure}[H]
\centering
\includegraphics[width=0.9\textwidth,height=0.9\textheight,keepaspectratio]{\fbpath/08.01 СВ/661.0121 ТЗНП/img/OA_TNZNP_without_T.pdf}
\captionof{figure}{Функциональная схема ОУ2 ТЗНП}\label{pict_tnznp:OA_TNZNP_without_T}
\end{figure}

\item
Оперативное ускорение ТЗНП со второй выдержкой времени предназначено для оперативного ускорения
ступени ТЗНП при выполнении переключений в сети. Ввод в работу ОУ2 ТЗНП и выбор ускоряемой
ступени задается программным переключателем \textbf{«Ступень ОУ2 ТЗНП»}. Предусмотрена возможность вывода
направленности ускоряемой ступени с помощью программного ключа \textbf{«Вывод напр при ср ТЗНП»}. Действие
ускорения осуществляется с регулируемой независимой выдержкой времени \textbf{«Тср ОУ2 ТЗНП»}. При
выводе направленности ускоряемой ступени ее блокировка от БНТ осуществляется с помощью программного
ключа \textbf{«БНТ ОУ2 ТЗНП»}.

\item
Оперативный ввод ОУ2 ТЗНП производится с помощью команды управления \textit{«ТЗНП/ОУ2 ТЗНП: Ввод»}.

\end{enumerate}


%%%%%%%%%%%%%%%%%%%%%%%%%%%%%%%%%%%%%%%%%%%%%%%%%% АУ %%%%%%%%%%%%%%%%%%%%%%%%%%%%%%%%%%%%%%%%%%%%%%%%%%
\item Автоматическое ускорение (АУ)

\begin{enumerate}

\item
В составе логики токовой направленной защиты нулевой последовательности предусмотрен пуск
автоматического ускорения ступеней ТЗНП.

\vspace{2mm}
\begin{figure}[H]
\centering
\includegraphics[width=0.9\textwidth,height=0.9\textheight,keepaspectratio]{\fbpath/08.01 СВ/661.0121 ТЗНП/img/AA_TNZNP.pdf}
\captionof{figure}{Функциональная схема пуска АУ ТЗНП}\label{pict_tnznp:AA_TNZNP}
\end{figure}

\item
Ввод в работу пуска АУ ТЗНП и выбор ускоряемой ступени задается программным переключателем
\textbf{«Ступень АУ ТЗНП»}. Предусмотрена возможность вывода направленности ускоряемой ступени с помощью
программного ключа \textbf{«Вывод напр при ср ТЗНП»}. При выводе направленности ускоряемой ступени ее
блокировка от БНТ осуществляется с помощью программного ключа \textbf{«БНТ АУ ТЗНП»}.

\end{enumerate}

\end{enumerate}